\documentclass[conference]{IEEEtran}
\usepackage{amsmath}
\usepackage{booktabs}
\usepackage{pgfplots}
\pgfplotsset{compat=1.17}

\begin{document}

\title{Parallel and Distributed K-Nearest-Neighbor Classification: Comparing POSIX Threads, OpenMP, and MPI}

% TODO: replace with your name and confirm the email format.
\author{
\IEEEauthorblockN{Your Name}
\IEEEauthorblockA{Department of Computer Science\\
Virginia Commonwealth University, VCU\\
Richmond, USA\\
yourname@vcu.edu}
}

\maketitle

\begin{abstract}
This report presents a serial baseline and three parallel implementations
of the k-nearest-neighbor (KNN) classifier: a shared-memory version using
POSIX threads, a shared-memory version using OpenMP, and a distributed
version using MPI. All three parallel versions partition the classification
work along the same axis, the set of test instances, so that every worker
reproduces the exact predictions of the serial version regardless of how
many workers are used. The implementations were developed and correctness
tested on \texttt{maple.cs.vcu.edu} and evaluated for runtime and speedup on
both \texttt{maple.cs.vcu.edu} and the \texttt{athena.hprc.vcu.edu} Slurm
cluster, using three datasets of increasing size and worker counts of 1, 2,
4, 8, 16, 32, 64, and 128. % TODO: one-sentence headline result once you have final numbers.
\end{abstract}

\begin{IEEEkeywords}
HPDS, KNN, POSIX threads, OpenMP, MPI, parallel computing, speedup
\end{IEEEkeywords}

\section{Introduction}

The k-nearest-neighbor classifier assigns a class to a query instance by
finding the $k$ closest instances to it in a labeled training set and
taking a majority vote of their labels. Its accuracy is competitive with
far more complex classifiers on many problems, but its cost grows as
$O(\mathit{num\_test} \times \mathit{num\_train} \times
\mathit{num\_features})$, since every test instance must be compared
against every training instance across every feature. On datasets with
tens of thousands of instances this cost becomes the practical limit on
how quickly the classifier can be run, which makes KNN a useful vehicle for
comparing parallelization strategies: the algorithm itself is simple
enough that the effect of the parallelization strategy, rather than the
complexity of the algorithm, dominates the results.

This report implements and evaluates three parallel strategies for the
same classification workload. All three are built around the same
partitioning decision, splitting the test instances across workers rather
than the training instances, for a reason argued in Section~\ref{sec:partition}:
it is the only one of the two splits that reproduces the serial version's
predictions exactly, including on ties, without any extra work.

The rest of this report is organized as follows. Section~\ref{sec:methodology}
describes the serial algorithm, the shared partitioning strategy, the three
parallel implementations, and the experimental setup and metrics used to
evaluate them. Section~\ref{sec:results} presents the measured runtimes,
speedups, and accuracy. Section~\ref{sec:analysis} discusses the trends in
those results and compares the three approaches. Section~\ref{sec:conclusion}
concludes.

\section{Methodology}
\label{sec:methodology}

\subsection{Serial KNN Algorithm}

The provided serial implementation classifies each test instance
independently. For a single test instance, it computes the Euclidean
distance to every training instance, maintains a sorted list of the $k$
closest training instances seen so far, and after scanning the full
training set predicts the class that appears most often among those $k$
neighbors. A new training instance is inserted into the candidate list only
when its distance is strictly smaller than an existing candidate's
distance, so among instances at exactly equal distance, the one encountered
earlier in the scan keeps its place. This tie-breaking rule matters for
Section~\ref{sec:partition}.

\subsection{Work Partitioning Strategy}
\label{sec:partition}

All three parallel implementations divide the set of test instances into
contiguous blocks and assign one block to each worker (thread or MPI
process), which then scans the entire training set for each test instance
in its block, exactly as the serial version does. Worker $t$ of $P$ workers
is given the half-open range $[\lfloor tn/P \rfloor, \lfloor (t+1)n/P
\rfloor)$ over $n$ test instances. This splits the test instances into $P$
blocks whose sizes differ by at most one, and it degrades gracefully when
$P$ exceeds $n$: a worker with $t \geq n$ receives an empty range and
simply does no work, which occurs in this evaluation whenever the small
dataset (160 test instances) is run with more than 160 workers.

An alternative partitioning splits the training set instead, and has every
worker scan a shard of the training instances for the same query,
building partial candidate lists that are then merged. This exposes more
parallel work, since the training sets used here are an order of magnitude
larger than the test sets, but it changes the order in which training
instances are compared for a given query. Because the serial algorithm
breaks distance ties by scan order, merging partial candidate lists in a
different order than the original single-threaded scan can change which
neighbor is kept when two training instances are exactly equidistant from
a query, changing the resulting prediction. Splitting on the test set
instead preserves the original scan order inside every worker, so every
implementation in this report reproduces the serial predictions exactly.
This is confirmed empirically in Section~\ref{sec:results}.

Splitting on the test set is also already load balanced without any
runtime scheduling. Every test instance costs exactly
$\mathit{num\_train} \times \mathit{num\_attributes}$ operations to
classify, so an equal-sized static split already assigns equal work to
every worker, and none of the three implementations uses a dynamic
work-stealing or work-queue scheme.

\subsection{POSIX Threads Implementation}

The pthreads version packs every value a worker thread needs, the shared
read-only dataset pointers, the worker's assigned range of test instances,
and a pointer to the shared output array, into a struct, and allocates one
struct instance per thread so that no two threads ever read or write the
same struct while a thread might still be starting up. Each worker writes
predictions only into the slice of the output array given by its own
range, so no lock is needed on the output array, since the ranges assigned
to different threads never overlap.

Each worker also allocates its own private copy of the two scratch
buffers used during classification, the sorted candidate list and the
per-class vote counts, once when the thread starts, and reuses that
private copy across every test instance in its range. The serial version
reuses a single copy of these buffers across all test instances because
only one instance is ever being classified at a time; sharing that same
buffer between multiple threads would let two threads overwrite each
other's in-progress candidate list, so each thread is given its own
independent copy instead. The main thread waits for every worker thread to
finish with \texttt{pthread\_join} before reading the completed
predictions array, which acts as the barrier that guarantees every write
is visible before the result is used.

\subsection{OpenMP Implementation}

The OpenMP version expresses the same partitioning using
\texttt{\#pragma omp parallel} to start a team of threads, followed by a
separate \texttt{\#pragma omp for}, with no \texttt{schedule} clause, to
divide the test-instance loop across that team. The two pragmas are kept separate,
rather than using the combined \texttt{parallel for} form, so that each
thread's private scratch buffers can be allocated once, between the two
pragmas, instead of once per test instance.

Those scratch buffers are declared as ordinary local variables inside the
parallel region rather than passed through OpenMP's \texttt{private} or
\texttt{firstprivate} data-sharing clauses. For a pointer variable, both
clauses only copy the pointer's value, an address, not the buffer that
address points to; every thread would therefore end up pointing at the
same shared buffer, or in the case of \texttt{private}, at an
uninitialized address. Declaring the buffers inside the parallel region
instead relies on the fact that every thread executing that region runs
the declaration itself and so receives its own independent storage,
without needing any explicit data-sharing clause. No \texttt{schedule}
clause is written because, as argued in Section~\ref{sec:partition}, every
test instance costs the same amount of work, so the fixed, contiguous
split every mainstream OpenMP runtime defaults to is already balanced, and
an explicit dynamic schedule would only add scheduling overhead for no
benefit here. The implicit barrier at the end of the \texttt{omp for} construct guarantees
every thread has finished writing its predictions before its scratch
buffers are freed and before the parallel region closes.

\subsection{MPI Implementation}

Unlike the two shared-memory versions, MPI ranks are separate processes
with separate address spaces, so data cannot be shared through a common
pointer. The implementation avoids communicating the dataset itself by
having every rank independently parse both the training and test files
before the timed region begins, so every rank ends up with its own
complete copy of the data without any dataset bytes crossing between
processes. This trades per-rank memory for simplicity and for eliminating
one whole category of communication, at the cost of not scaling to
datasets too large to fit in a single node's memory, which is not a
limitation for the datasets used here.

Every rank computes the same partition table used by the threaded version,
covering every rank's range, not only its own, since the collection step
described next needs the full table. Each rank then classifies its own
range of test instances into a local array sized to exactly its own share
of the work. Because each rank's scratch buffers and local results array
exist only in that rank's own memory, there is no data race to avoid here
the way there is with threads.

The one step with no equivalent in the threaded versions is collecting
every rank's local results back into a single array on rank~0, which alone
computes the confusion matrix. This is done with \texttt{MPI\_Gatherv}
rather than the simpler \texttt{MPI\_Gather}, because \texttt{MPI\_Gather}
requires every rank to contribute the same number of elements, while the
partition table gives blocks that differ by one whenever the test set does
not divide evenly across ranks. \texttt{MPI\_Gatherv} takes a per-rank
count and a per-rank displacement, so each rank's contribution lands at
the correct global offset regardless of block size, and rank~0 recovers
the predictions in their original order. \texttt{MPI\_Gatherv} is also a
synchronization point, so an \texttt{MPI\_Barrier} is placed immediately
before the timer starts, to prevent a rank that happens to finish parsing
its input files earlier than the others from including that head start,
and the resulting idle wait inside \texttt{MPI\_Gatherv}, in its own
measured runtime.

\subsection{Experimental Setup}

% TODO: confirm CPU count / memory for maple and the node count actually used
% on Athena for your runs, and adjust this paragraph if it differs.
Development, debugging, and correctness testing were done on
\texttt{maple.cs.vcu.edu}, a single node with 28 logical CPUs. Full timing
experiments were run on both \texttt{maple.cs.vcu.edu} and the
\texttt{athena.hprc.vcu.edu} Slurm cluster. On Athena, the MPI
implementation was run twice for every configuration, once with every rank
forced onto a single node (\texttt{--nodes 1}), and once allowed to spread
across multiple nodes, to compare the cost of intra-node versus
inter-node communication for the same workload.

Three datasets of increasing size were used, summarized in
Table~\ref{tab:datasets}. Each was evaluated with $k=3$ and with worker
counts of 1, 2, 4, 8, 16, 32, 64, and 128 threads or processes. Every
configuration was run three times and the runtimes reported in
Section~\ref{sec:results} are the mean of those three runs.

\begin{table}[htbp]
\caption{Dataset sizes}
\label{tab:datasets}
\centering
\begin{tabular}{lrrr}
\toprule
Dataset & Train instances & Test instances & Attributes \\
\midrule
small  & 1{,}184  & 160   & 7 \\
medium & 14{,}708 & 740   & 11 \\
large  & 61{,}606 & 3{,}436 & 11 \\
\bottomrule
\end{tabular}
\end{table}

\subsection{Performance Metrics}

Two speedup definitions are reported, since they answer different
questions. Absolute speedup compares a parallel run against the serial
program,
\begin{equation}
S_{\text{abs}}(p) = \frac{T_{\text{serial}}}{T_{\text{parallel}}(p)},
\label{eq:abs}
\end{equation}
and reflects the real-world benefit of parallelizing, including whatever
fixed overhead the parallel version pays. Relative speedup instead
compares a parallel implementation against its own single-worker run,
\begin{equation}
S_{\text{rel}}(p) = \frac{T_{\text{parallel}}(1)}{T_{\text{parallel}}(p)},
\label{eq:rel}
\end{equation}
which isolates how well that implementation scales once its own fixed
overhead is factored out. Absolute speedup is used as the headline metric
in Section~\ref{sec:results}, with relative speedup reported alongside it.

The parallel fraction of the code is estimated from the absolute speedup
using the Karp-Flatt metric,
\begin{equation}
e(p) = \frac{1/S_{\text{abs}}(p) - 1/p}{1 - 1/p}, \qquad
f_{\text{parallel}}(p) = 1 - e(p),
\label{eq:karpflatt}
\end{equation}
where $e(p)$ is the estimated serial fraction at $p$ workers. Unlike
Amdahl's law solved for a single fixed serial fraction, Karp-Flatt is
evaluated at every measured value of $p$; an estimated serial fraction that
grows as $p$ increases is itself evidence of parallel overhead, such as
load imbalance or communication cost, that a model with one fixed serial
fraction cannot capture.

\section{Results}
\label{sec:results}

\subsection{Accuracy}

Table~\ref{tab:accuracy} reports the classification accuracy of the serial
implementation on each dataset. Every parallel implementation, at every
worker count from 1 to 128, reproduced these accuracies exactly, confirmed
with an automated check comparing each parallel run's reported accuracy
against the serial run's accuracy for the same dataset. The lower accuracy
on the medium dataset is a property of that dataset's class separability,
not a defect in any implementation: it is reproduced identically by the
serial version and by every parallel version.

\begin{table}[htbp]
\caption{Serial classification accuracy, $k=3$}
\label{tab:accuracy}
\centering
\begin{tabular}{lr}
\toprule
Dataset & Accuracy \\
\midrule
small  & 85.00\% \\
medium & 27.57\% \\
large  & 99.48\% \\
\bottomrule
\end{tabular}
\end{table}

\subsection{Runtime and Speedup on Maple}

% TODO: replace the placeholder rows below with the real values from
% results/summary.csv after running run_experiments.sh and analyze_results.py
% on maple with an unrestricted core allocation.
Table~\ref{tab:maple-threaded} and Table~\ref{tab:maple-openmp} report mean
runtime, absolute speedup, relative speedup, and estimated parallel
fraction for the pthreads and OpenMP implementations on maple, for the
large dataset. Figure~\ref{fig:maple-speedup} plots absolute speedup
against worker count for both implementations across all three datasets.

\begin{table}[htbp]
\caption{Pthreads on maple, large dataset (placeholder values)}
\label{tab:maple-threaded}
\centering
\begin{tabular}{rrrrr}
\toprule
Workers & Time (ms) & $S_{\text{abs}}$ & $S_{\text{rel}}$ & $f_{\text{parallel}}$ \\
\midrule
1   & TODO & TODO & 1.00 & --- \\
2   & TODO & TODO & TODO & TODO \\
4   & TODO & TODO & TODO & TODO \\
8   & TODO & TODO & TODO & TODO \\
16  & TODO & TODO & TODO & TODO \\
32  & TODO & TODO & TODO & TODO \\
64  & TODO & TODO & TODO & TODO \\
128 & TODO & TODO & TODO & TODO \\
\bottomrule
\end{tabular}
\end{table}

\begin{table}[htbp]
\caption{OpenMP on maple, large dataset (placeholder values)}
\label{tab:maple-openmp}
\centering
\begin{tabular}{rrrrr}
\toprule
Workers & Time (ms) & $S_{\text{abs}}$ & $S_{\text{rel}}$ & $f_{\text{parallel}}$ \\
\midrule
1   & TODO & TODO & 1.00 & --- \\
2   & TODO & TODO & TODO & TODO \\
4   & TODO & TODO & TODO & TODO \\
8   & TODO & TODO & TODO & TODO \\
16  & TODO & TODO & TODO & TODO \\
32  & TODO & TODO & TODO & TODO \\
64  & TODO & TODO & TODO & TODO \\
128 & TODO & TODO & TODO & TODO \\
\bottomrule
\end{tabular}
\end{table}

\begin{figure}[htbp]
\centering
\begin{tikzpicture}
\begin{axis}[
    width=\linewidth, height=5.5cm,
    xlabel={Workers}, ylabel={Absolute speedup},
    xmode=log, log basis x=2,
    legend pos=north west, legend style={font=\scriptsize},
]
\addplot table[x=workers, y=speedup_abs] {data/maple_threaded_large.dat};
\addplot table[x=workers, y=speedup_abs] {data/maple_openmp_large.dat};
\legend{Pthreads, OpenMP}
\end{axis}
\end{tikzpicture}
\caption{Absolute speedup vs. worker count on maple, large dataset
(placeholder data; replace \texttt{data/*.dat} with the files written by
\texttt{analyze\_results.py} before compiling the final report).}
\label{fig:maple-speedup}
\end{figure}

\subsection{Runtime and Speedup on Athena}

% TODO: replace the placeholder rows below with the real values from your
% Slurm runs.
Table~\ref{tab:athena-mpi} reports the MPI results on Athena for the large
dataset, comparing ranks forced onto a single node against ranks allowed to
spread across multiple nodes. Figure~\ref{fig:athena-mpi} plots absolute
speedup for both placement policies across all three datasets.

\begin{table}[htbp]
\caption{MPI on Athena, large dataset, single node vs. multiple nodes (placeholder values)}
\label{tab:athena-mpi}
\centering
\begin{tabular}{rrrrr}
\toprule
 & \multicolumn{2}{c}{Forced 1 node} & \multicolumn{2}{c}{Any number of nodes} \\
Ranks & Time (ms) & $S_{\text{abs}}$ & Time (ms) & $S_{\text{abs}}$ \\
\midrule
1   & TODO & TODO & TODO & TODO \\
2   & TODO & TODO & TODO & TODO \\
4   & TODO & TODO & TODO & TODO \\
8   & TODO & TODO & TODO & TODO \\
16  & TODO & TODO & TODO & TODO \\
32  & TODO & TODO & TODO & TODO \\
64  & TODO & TODO & TODO & TODO \\
128 & TODO & TODO & TODO & TODO \\
\bottomrule
\end{tabular}
\end{table}

\begin{figure}[htbp]
\centering
\begin{tikzpicture}
\begin{axis}[
    width=\linewidth, height=5.5cm,
    xlabel={MPI ranks}, ylabel={Absolute speedup},
    xmode=log, log basis x=2,
    legend pos=north west, legend style={font=\scriptsize},
]
\addplot table[x=workers, y=speedup_abs] {data/athena_mpi_forced_1_node_large.dat};
\addplot table[x=workers, y=speedup_abs] {data/athena_mpi_any_nodes_large.dat};
\legend{Forced 1 node, Any number of nodes}
\end{axis}
\end{tikzpicture}
\caption{Absolute speedup vs. rank count on Athena, large dataset,
single-node vs. multi-node placement (placeholder data).}
\label{fig:athena-mpi}
\end{figure}

\section{Analysis and Discussion}
\label{sec:analysis}

% TODO: this section argues the shape of trend to expect and why; once real
% numbers are in Section III, replace the conditional phrasing ("if speedup
% plateaus...") with a direct statement of what was actually observed, and
% adjust any claim the real data contradicts.

Because every test instance costs the same amount of work, the pthreads
and OpenMP versions are not expected to show meaningful imbalance between
workers at any worker count. Where their speedup curves are expected to
depart from linear is at worker counts beyond the number of physical cores
available, since maple's 28 logical CPUs place a hard ceiling on how many
computations can genuinely happen at once regardless of how many threads
are requested. If the absolute-speedup curve flattens or turns downward
before reaching the highest tested worker counts, that is expected
behavior once thread count exceeds available cores, not a defect in the
partitioning scheme, and the Karp-Flatt estimated parallel fraction should
be read as capturing that overhead rather than as a fixed property of the
algorithm.

The small dataset, with only 160 test instances, is expected to show this
ceiling earliest and most sharply, since at 128 workers each worker
receives at most two test instances, and the fixed cost of creating a
thread or an MPI rank is no longer small next to the amount of work that
thread or rank actually performs. This is the same reasoning discussed in
the accompanying methodology explanations for why the partitioning formula
must handle worker counts larger than the number of test instances
gracefully.

For MPI, the comparison between ranks forced onto one node and ranks
allowed to spread across nodes isolates the cost of the single
communication step in this implementation, the \texttt{MPI\_Gatherv} call
that assembles the final predictions on rank~0. Intra-node communication
goes through shared memory, while inter-node communication goes over the
cluster interconnect; if the multi-node curve in
Figure~\ref{fig:athena-mpi} sits visibly below the single-node curve at
higher rank counts, that gap is attributable to this difference in
communication cost, since the two runs execute identical computation and
differ only in where their ranks are physically placed.

% TODO: add a direct three-way comparison paragraph here once all data is in,
% covering which implementation reached the highest speedup and why (e.g.
% MPI's per-call communication overhead vs. threads' shared-memory writes).

\section{Conclusion}
\label{sec:conclusion}

This report implemented and evaluated three parallel strategies, POSIX
threads, OpenMP, and MPI, for the k-nearest-neighbor classifier, all built
around the same test-instance partitioning strategy. That shared strategy
was chosen specifically because it reproduces the serial version's
predictions exactly regardless of worker count, which was confirmed
empirically across all three implementations, three datasets, and every
tested worker count.
% TODO: one or two closing sentences once the real numbers support a
% specific claim, e.g. which implementation scaled best on which dataset
% size and why, referring back to Section IV.

\begin{thebibliography}{9}
\bibitem{knn} T. Cover and P. Hart, ``Nearest neighbor pattern classification,'' \emph{IEEE Transactions on Information Theory}, vol. 13, no. 1, pp. 21--27, 1967.
\bibitem{pthreads} B. Lewis and D. J. Berg, \emph{Multithreaded Programming with Pthreads}. Prentice-Hall, Inc., 1998.
\bibitem{openmp} OpenMP Architecture Review Board, ``OpenMP Application Programming Interface, Version 5.2,'' 2021.
\bibitem{mpi} W. Gropp, E. Lusk, and A. Skjellum, \emph{Using MPI: Portable Parallel Programming with the Message-Passing Interface}, 3rd ed. MIT Press, 2014.
\bibitem{karpflatt} A. H. Karp and H. P. Flatt, ``Measuring parallel processor performance,'' \emph{Communications of the ACM}, vol. 33, no. 5, pp. 539--543, 1990.
\end{thebibliography}

\end{document}
